\section{Intended use cases}

The resource supports five classes of use.

\begin{enumerate}
  \item \textbf{Independent qualitative replication.} A coder can analyze the blind hard-case set from the linked public sources and compare classifications only after submission.
  \item \textbf{Cross-community evidence studies.} Researchers can apply the unit, provenance, uncertainty, and claim-boundary pattern to another specialist technical community and compare which constructs transfer.
  \item \textbf{Laboratory knowledge stewardship.} Teams can adapt the question and resolved-knowledge schemas for internal or public support systems while retaining applicability, evidence, ownership, and expiry.
  \item \textbf{Operational instrumentation.} Laboratories can use the interface and event schemas to collect real integration and incident data with denominators unavailable in forum evidence.
  \item \textbf{AI-method evaluation.} Benchmark designers can construct missing-context tasks and report performance separately across syntactic, simulation, dry, wet, assay, and production stages.
\end{enumerate}

The resource should be evaluated by outcomes appropriate to each use. Coding replication requires agreement and adjudication records. Knowledge interventions require clarification burden, resolution time, recurrence, staleness, and correction latency. Interface registries require integration time, executable-device coverage, and failed dispatches. Event models require diagnosis time, safe-state time, salvage, final disposition, and recurrence. AI benchmarks require clarification quality, unsafe omission, human correction, and validation-stage performance.

The resource should not be used to infer vendor reliability, market position, installed base, contributor competence, or field-wide failure incidence. Those interpretations lack the required sampling frame and operational denominator.
